% docs/manual/operations/chapters/04-performance.tex
% 第 4 章：性能调优

\chapter{性能调优}

性能调优是运维的核心问题：什么时候 \openbench{} 跑得慢，是 CPU 在等 IO，还是内存在 swap，还是并行根本没用上？本章分类讲解三大瓶颈与对应的诊断/调优手段。

\section{三大瓶颈}

任何长跑任务的性能问题都落在三类之一：

\begin{longtable}{p{2cm} p{4.5cm} p{6cm}}
  \toprule
  瓶颈 & 症状 & 工具 \\
  \midrule
  \endhead
  CPU & 全核高负载，进度仍慢 & \texttt{top -H}、\texttt{htop}、py-spy \\
  内存 & swap、OOM kill、运行越久越慢 & \texttt{free}、\texttt{/usr/bin/time -v}、\texttt{psutil} \\
  I/O & CPU < 30\%、disk wait 高 & \texttt{iostat -x}、\texttt{iotop}、NFS 上 \texttt{nfsstat} \\
  \bottomrule
\end{longtable}

先用 \texttt{/usr/bin/time -v openbench run config.yaml} 跑一次，看 \texttt{Elapsed (wall clock)}、\texttt{Maximum resident set size}、\texttt{Percent of CPU this job got}：CPU\% < 50\% 说明在等 IO 或并行不够。

\section{num\_cores 选择}

\subsection{自动检测 vs 显式指定}

\begin{minted}{yaml}
project:
  num_cores: null    # 默认：os.cpu_count()
  num_cores: 16      # 显式（推荐）
\end{minted}

显式优于自动的原因：

\begin{itemize}
  \item HPC 上 \texttt{os.cpu_count()} 看到的是节点全部核（如 64），但 SLURM 只分配你 16 核 → oversubscription
  \item 笔记本上 hyperthreading 让 \texttt{cpu\_count} 翻倍 → 实际 8 物理核被当 16 用
  \item 显式数字让性能可复现
\end{itemize}

\subsection{与 SLURM 对齐}

\begin{minted}{bash}
openbench run config.yaml --cores $SLURM_CPUS_PER_TASK
\end{minted}

或在 yaml 里：

\begin{minted}{yaml}
project:
  num_cores: 32   # 与 #SBATCH --cpus-per-task=32 一致
\end{minted}

\subsection{选多少最佳}

\begin{itemize}
  \item I/O 密集任务（站点评估）：核数 = 物理 CPU / 2（多了 IO 拥塞抵消并行收益）
  \item CPU 密集任务（统计模块）：核数 = 物理 CPU
  \item 内存紧张：核数 = 总内存 / 单 task 峰值（避免 OOM）
\end{itemize}

\section{当前并行边界}

\openbench{} v3.0a1 的并行实际作用范围：

\subsection{已并行}

\begin{itemize}
  \item \textbf{站点级}：\modname{data.processing.DatasetProcessing} 用 \modname{joblib} 并行处理多个站点
  \item \textbf{按年合并}：\modname{data.pipeline} 在合并多年 NC 时按年并发读
  \item \textbf{xarray 内部 dask}：metric 计算时部分 chunked 操作走 dask
\end{itemize}

\subsection{尚未并行}

\begin{itemize}
  \item \textbf{Task 级（变量级）}：多个 (var × sim × ref) task 仍串行执行；reserved
  \item \textbf{visualization 阶段}：图序列化生成
  \item \textbf{Statistics 阶段}：模块串行
\end{itemize}

\warninline{这意味着 \yamlkey{num_cores: 64} 在 grid 评估上不一定比 \yamlkey{num_cores: 16} 快很多——除非你的评估有大量站点（站点级并行才放大）。task 级并行在 Plan 5 之后的版本会启用。}

\section{内存监控}

\subsection{util.memory 的工具}

\modname{util.memory.cleanup\_memory()}：

\begin{itemize}
  \item Python 缓存（functools / importlib / linecache）
  \item 科学库缓存（numpy / xarray / pandas / matplotlib）
  \item \texttt{gc.collect()} 全代
  \item 弱引用清理
\end{itemize}

主流程在 task 之间会调用它。手动用：

\begin{minted}{python}
from openbench.util.memory import cleanup_memory
cleanup_memory(verbose=True)
\end{minted}

\subsection{psutil 监控}

\begin{minted}{python}
import psutil
proc = psutil.Process()
print(f"RSS: {proc.memory_info().rss / 1e9:.2f} GB")
\end{minted}

\subsection{峰值内存}

\begin{minted}{bash}
/usr/bin/time -v openbench run config.yaml 2>&1 | grep "Maximum resident"
# Maximum resident set size (kbytes): 18234560
\end{minted}

\subsection{OOM 三种情形}

\textbf{情形 A：全球 hourly + 多年 + 多 sim}：
\begin{itemize}
  \item 单文件可能 30 GB+（hourly × 5yr × 0.25°）
  \item \openbench{} 默认 lazy load，但某些 metric 强制 load 全部
  \item 缩 \yamlkey{years} 到 1 年试试
\end{itemize}

\textbf{情形 B：dask chunk 太大}：
\begin{itemize}
  \item \openbench{} 不直接配 dask，但 xarray 调用走 dask 默认配置
  \item 看 \texttt{xr.open\_dataset(chunks=\{...\})} 是否合理（grid evaluator 用默认）
  \item 可在调用前设 \texttt{xr.set\_options(chunk='auto')}
\end{itemize}

\textbf{情形 C：站点 CSV + grid 笛卡尔积}：
\begin{itemize}
  \item 千站点 × 全球 grid → 中间内存暴增
  \item station\_matcher 应该按站点 streaming 处理（已实现），但某些路径会 materialize
  \item 缩区域或减少站点列表
\end{itemize}

\section{I/O 优化}

\subsection{HDF5 文件锁}

\openbench{} 已自动设：

\begin{minted}{python}
os.environ.setdefault("HDF5_USE_FILE_LOCKING", "FALSE")
\end{minted}

\warninline{在 Notebook / 自定义脚本里手动调 \modname{run\_evaluation} 时，必须 \emph{在 import 前} 设这个环境变量，否则不生效。}

\subsection{NFS / Lustre 上的小文件抖动}

reference 数据集动辄数千个 NC 文件，每个 ~MB。NFS 上每次 \texttt{open()} 都走元数据服务器。

\textbf{对策}：

\begin{itemize}
  \item 让 reference 数据放在 Lustre（HPC 标配）而非家目录 NFS
  \item 启用 \texttt{open\_mfdataset} 一次打开多文件，减少 round trips
  \item 把多年小文件预合并为一个大 NC（看 \cli{cdo cat *.nc merged.nc}）
\end{itemize}

\subsection{zarr cache}

\openbench{} 默认在 \file{output/<run>/data/} 写 zarr 中间产物。第二次跑同 task 直接读 zarr，避开 NC 解析。

\textbf{陷阱}：

\begin{itemize}
  \item 第一次跑加上 zarr 写入时间；总时间可能更长
  \item 多个评估目录都会各自有一份 cache → 占空间。多用户共享时考虑统一 cache 目录
\end{itemize}

\section{Dask scheduler 选择}

\modname{xarray} 的某些算子走 \modname{dask}。\openbench{} 不直接配 dask，但可通过环境变量影响：

\begin{minted}{bash}
export DASK_SCHEDULER="threads"     # 默认；适合多数评估
export DASK_SCHEDULER="processes"   # 跨进程；某些 IO bound 场景更快
\end{minted}

或在 Python 里：

\begin{minted}{python}
import dask
dask.config.set(scheduler='threads')
\end{minted}

\section{Oversubscription 风险}

嵌套并行陷阱：joblib（核数 N）× MKL（核数 M）= N × M 线程，远超物理核数。

\openbench{} 的处置：

\begin{itemize}
  \item 启动时设 \texttt{OMP\_NUM\_THREADS=1}（如果还没设）
  \item joblib 用 \texttt{loky} backend，子进程独立 BLAS
\end{itemize}

\textbf{用户可做}：

\begin{minted}{bash}
export OMP_NUM_THREADS=1
export MKL_NUM_THREADS=1
export OPENBLAS_NUM_THREADS=1
openbench run config.yaml --cores 16
\end{minted}

\section{实测建议}

\subsection{先小规模 profiling}

\begin{minted}{yaml}
# 先跑 1 个变量 + 1 sim + 1 年
project:
  years: [2010, 2010]   # 1 年
evaluation:
  variables: [Evapotranspiration]   # 1 变量
simulation:
  test:
    model: CoLM2024
    root_dir: ...
\end{minted}

\begin{minted}{bash}
/usr/bin/time -v openbench run small.yaml --cores 4 2>&1 | tee small.log
\end{minted}

得出单 task 时间与内存，再外推到完整评估。

\subsection{瓶颈定位顺序}

\begin{enumerate}
  \item 看 \texttt{Percent of CPU this job got}
    \begin{itemize}
      \item < 50\% → IO 或并行不够
      \item > 90\% → CPU bound
    \end{itemize}
  \item 看 \texttt{Maximum resident set size}
    \begin{itemize}
      \item > 80\% 物理内存 → 内存压力大
      \item < 40\% 物理内存 → 可以增 \yamlkey{num_cores}
    \end{itemize}
  \item \texttt{iostat -x 1} 看 \%util
    \begin{itemize}
      \item disk \%util > 80\% → IO bound
    \end{itemize}
\end{enumerate}

\section{调优 checklist}

升级到生产前过一遍：

\begin{itemize}
  \item \yamlkey{num_cores} 显式设为合理值（不是默认 null）
  \item 设 \texttt{OMP\_NUM\_THREADS=1} 防嵌套并行
  \item Reference 数据放高速盘（不在家目录 NFS）
  \item 第一次小规模 profile，记下内存/时间
  \item 跑 \cli{check} 通过且无 warning
  \item 用 SLURM 提交（不要在 login 节点直接跑长任务）
  \item \yamlkey{generate_report: true} 不影响主流程速度太多但加 5-15\%；要赶时间可关
\end{itemize}

\section{下一步}

下一章详解缓存策略：增量缓存、清理、多用户共享、损坏检测。
